\chapter{Experiment Provenance and Reproducibility}
\label{app:reproducibility}

\section{Available Experiment Records}
The experiment records identify the versions of datasets, model code and configurations
used in each study. They include selected predictions, training records, completion
receipts and result summaries. Some raw data, trained model weights, feature caches and
historical prediction files are held separately. The source manifests record their
locations where known.

The numerical checks used existing outputs. The original training, gradient checks and
saved-model inference belong to the recorded experiments. The checks described below did
not involve refitting a model or running new inference on raw inputs.\ev{External
inspection scope}{\EvidenceRef{E106}{}}

\section{Numerical Verification}
For REASSEMBLE, average precision was recalculated from saved predictions for the core
and specified follow-up comparisons. For MSM, the check used the primary assessment
probabilities for intact inputs and selected perturbations. For the original ViFailback
experiment it used episode-level branch and fusion probabilities. For HANDOVER, the five
clean means and seed standard deviations were recalculated from three sets of per-seed
metrics. The retained export did not include individual trial probabilities.
The PHM comparison used the original run summary and paired-bootstrap record.

\section{Reconstructing an Evaluation Cohort}
The REASSEMBLE registry links each evaluation to the native recording and segment keys
of its assessed actions. It distinguishes training, validation, calibration and
assessment membership. Some legacy identifiers are aliases. The recorded fit identity
also distinguishes runs that share an outer seed label but use different random states
for model training.

The 753 actions in the selected outer fold form a subset of the pooled 4,530-action
assessment. Native recording and segment identifiers establish common test membership
across models on the smaller set. Checkpoint selection, calibration and threshold
procedures differed between some fits on that set, limiting attribution of a score
difference to architecture.\ev{Cohort reconstruction instructions}{\EvidenceRef{E050}{}}

MSM links intervals by their native IDs, label~2 policy and experiment or date group.
HANDOVER links trial IDs within the official split and pairs video with the common
wrench window. Participant identities are unresolved. ViFailback uses the recorded
original and converted release revisions, canonical task identifiers and the final
cohort with three-camera quality checks. Its original and matched studies share 4,382
episodes but use different fitting procedures.

\section{Hashes and Source-State Identity}
The external-source manifest gives SHA-256 digests for files at their original location
and for local copies. It also identifies the Git revision used for the checks. Some reports
included uncommitted changes, so the commit identifier alone cannot reconstruct those
files.

The recorded original and local hashes can differ after line-ending conversion. The
source path and corresponding hash identify which version underlies a calculation.
ViFailback completion records identify the condition, training seed and fit role for
each completed comparison in Table~\ref{tab:vif-status}.\ev{Source-copy
identities}{\EvidenceRef{E107}{}}

\section{Saved-Result Checking and Experiment Replay}
The archive's integrity checks compare file digests and references in the cohort
registry. Score checks match the assessment identities, labels and probability
definitions before computing performance. Failure-positive average precision uses
precision at distinct score thresholds and increments in recall. Trapezoidal area under
the precision--recall curve uses a different interpolation.

Pooling disjoint outer-fold predictions before computing AP weights the resulting joint
ranking. Averaging AP over folds gives each fold the specified aggregation weight.
Similarly, averaging probabilities across seeds defines an ensemble, whereas averaging
AP across seeds describes variation among fitted predictors. The uncertainty estimates
use the study-specific cluster units. In HANDOVER the absence of participant identities
prevents a participant-cluster interval.

A training replication also depends on the original raw inputs, model and configuration
versions, software dependencies and either the fitted weights or their initialisation
procedure, depending on the stage being repeated. Where only a summary survives, it is
possible to check the reported aggregation but not to regenerate individual predictions
from that summary.

\section{Coverage of the Checks}
The matched two-epoch REASSEMBLE comparison records new and reused fits and selected
saved-model checks on pixel inputs. Those checks leave the cause of its seed-dependent
behaviour unresolved. In ViFailback, FROZEN20 and DONOR20 provide complete three-seed
assessments, FULL20 provides one outer-fold comparison, and the three hundred-epoch
conditions provide separate single-role pilots. The remaining FULL20 fits and SQ3
corruption tests have no measured results. The completed records support checking saved
results within those scopes; they do not constitute independent training replication.\ev{Diagnostic replay
scope}{\EvidenceRef{E060}{}, \EvidenceRef{E113}{}, \EvidenceRef{E114}{}}

\section{Verification of Reported Metrics}
The retained verification recalculated failure-positive AP, Brier score and fifteen-bin
ECE from the retained core REASSEMBLE arrays for clean inputs, all seven strongest
corruptions and both known-absence conditions. All 32 model rows matched the saved
summary within $6.2\times10^{-16}$; labels and fold vectors matched across conditions.
This was a saved-score calculation, without new model fitting or inference, and did not
repeat bootstrap intervals or independently reconstruct recording identities.
For MSM, all sixty clean summary means matched the means of their three saved pooled
seed APs; the visual in-sample diagnostic averaged nine saved fit scores.
These checks verified aggregation of existing scores rather than regenerating predictions.
PHM values were checked against retained metric tables and reports, with historical
source gaps preserved. ViFailback calibration coefficients and control scores remain
attributed to the retained reports: individual matched prediction/calibrator files were
not locally available for a new reconstruction. Earlier checks
retain their own dates and access scope.\ev{Recorded verification of 27 November 2026 and its inspected inputs}
{\EvidenceRef{E133}{}, \EvidenceRef{E135}{}, \EvidenceRef{E137}{}, \EvidenceRef{E139}{}}

\section{Project Evidence File Register}
\label{sec:evidence-files}
Each identifier refers to one source file or dated copy. The \enquote{Link here}
entry links directly to that file at a verified GitHub commit. The accompanying file
index gives source paths, recorded hashes, publication revisions and the sections,
code routines or samples cited. New entries prefixed \texttt{repo:} in that
index resolve from the research-repository root. For E142--E152, the register itself
gives the repository-relative paths, and the accompanying provenance records retain
the checked SHA-256 identities. The preserved package contains
baseline hashes, the source diff and retained saved-score checks. Its ViFailback check
verified 1,008 scalar metrics from 168 saved model/cohort vectors and the completed-fit
records. The paired intervals remain those recorded on 24 November; the saved-score
check did not repeat their bootstrap calculation.\ev{Saved-score verification recorded on 25 November 2026}{\EvidenceRef{E111}{}, \EvidenceRef{E121}{}}

E093 links to a hash-checked archived copy recovered from the thesis evidence package;
the original file location is no longer available. E152 identifies the retained final snapshot of 24 November 2026,
which includes the FULL100 pilot, 22 completed outer fusion assessments and six pooled
summaries. These are separate dated versions.

\begingroup
\singlespacing\small
\setlength{\tabcolsep}{3pt}
\setlength{\LTleft}{0pt}
\setlength{\LTright}{\fill}
\captionsetup{hypcap=false}
\captionof{table}{Project Evidence Files and Repository Links}\label{tab:project-evidence}
\begin{longtable*}{@{}P{0.075\linewidth}P{0.19\linewidth}P{0.475\linewidth}P{0.21\linewidth}@{}}
\toprule ID & Study or source & File & Repository link\\\midrule\endfirsthead
\toprule ID & Study or source & File & Repository link\\\midrule\endhead
\bottomrule\endfoot
\PrintEvidenceRows
\end{longtable*}
\endgroup
